% !TeX root = ../教材题库_全部条目_转写工作稿.tex
% Source: XJ-B1-C04-EX-P096-Y0978
% 第4章 / 4.1 实数指数幂和幂函数 / 教材 p.96 / PDF 第103页
% OCR未校对：公式与文字均可能有误；未完成数学TeX转写；禁止直接发给学生
% 原题及相邻原解截图（核对用）：../题目原图/XJ-B1-C04-EX-P096-Y0978-01.webp
% 本文件只在转写完成后作为题干使用；OCR 条目尚不是可靠题干。
\begin{bankitem}{4.1 例题 例题3}{第4章 / 4.1 实数指数幂和幂函数 / 教材 p.96 / PDF 第103页}{XJ-B1-C04-EX-P096-Y0978}{OCR未校对：公式与文字均可能有误；未完成数学TeX转写；禁止直接发给学生}
\begin{ocrdraft}
[教材p.96 / PDF第103页]
\par
0河用分数指数幂的形式表示下列各式（>0）：
\par
（2）a2
\par
（2）是
\par
（3）磊\ensuremath{\cdot}、\ensuremath{\text{\textperthousand}}一伍\ensuremath{\cdot}a3）2
\par
（a3）2
\par
a32
\end{ocrdraft}
\end{bankitem}
